\documentclass[../../../include/open-logic-section]{subfiles}

\begin{document}

\olfileid{sth}{story}{approach}
\olsection{સંચયી-પુનરાવર્તી અભિગમ}

ગણોને તબક્કાઓમાં કેવી રીતે ગોઠવીશું તેનું થોડું વધુ પૂર્ણ વર્ણન
આ છે:
\begin{quote}
ગણો \emph{તબક્કાઓ}માં રચાય છે. દરેક તબક્કા $S$ માટે કેટલાક
તબક્કાઓ $S$ની \emph{પહેલાં} હોય છે. તબક્કા $S$ પર, $S$ પહેલાંના
તબક્કાઓમાં રચાયેલા ગણોથી બનેલો દરેક સંગ્રહ ગણ તરીકે રચાય છે.
તબક્કાઓમાં રચાતા ગણો સિવાય બીજા કોઈ ગણો નથી.
\citep[p.~323]{Shoenfield:AST}
\end{quote}
આ \emph{ગણની સંચયી-પુનરાવર્તી કલ્પના}ની રૂપરેખા છે.
\olref[sth][][]{part}માં રજૂ કરવાના ઔપચારિક ગણસિદ્ધાંતનો તે
આધાર બનશે.

આને થોડું વધુ વિગતે તપાસીએ. શોએનફેલ્ડ પ્રક્રિયાનું વર્ણન કરે છે
તે મુજબ દરેક તબક્કે પહેલાંના તબક્કાઓમાંથી ઉપલબ્ધ થયેલા ગણો
વડે નવા ગણો રચીએ છીએ. તેથી શોએનફેલ્ડના વર્ણન મુજબ પ્રારંભિક
તબક્કા, તબક્કા $0$,ની \emph{પહેલાં} કોઈ તબક્કા નથી, અને
\emph{તો ખાસ કરીને} પહેલાંના તબક્કાઓમાંથી કોઈ ગણો ઉપલબ્ધ
નથી.\footnote{\emph{પ્રથમ} તબક્કો છે એવું શા માટે માનવું?
\olref[z][story]{sec}માં \stagesord{}ની ફૂટનોટ જુઓ.}
તેથી માત્ર એક ગણ રચીએ છીએ: કોઈ !!{element}s વિનાનો ગણ
$\emptyset$. તબક્કા $1$ પર પહેલાંના તબક્કાઓમાંથી બરાબર એક ગણ
ઉપલબ્ધ છે, તેથી માત્ર એક નવો ગણ $\{\emptyset\}$ રચાય છે.
તબક્કા $2$ પર પહેલાંના તબક્કાઓમાંથી બે ગણો ઉપલબ્ધ છે, અને
બે નવા ગણો $\{\{\emptyset\}\}$ અને
$\{\emptyset, \{\emptyset\}\}$ રચીએ છીએ. તબક્કા $3$ પર
પહેલાંના તબક્કાઓમાંથી ચાર ગણો ઉપલબ્ધ છે, તેથી બાર નવા ગણો
રચીએ છીએ\ldots. આમ ગણોનું સંચયી-પુનરાવર્તી ચિત્ર કંઈક નીચે
જેવું હશે (અંકો તબક્કાઓ દર્શાવે છે):
\begin{center}
\begin{tikzpicture}[scale=0.6]
	\tikzset{cut_here/.style={densely dotted}}
	\draw (-4,6) -- (-1, 0)-- (2,6); 
	\draw[cut_here] (-5,8)--(-4,6);
	\draw[cut_here] (2,6)--(3,8);
	\draw(-1.5, 1)--(-0.5, 1);
	\draw(-2, 2)--(0, 2);
	\draw(-2.5, 3)--(0.5,3);
	\draw(-3, 4)--(1, 4);
	\draw(-3.5, 5)--(1.5, 5);
	\draw(-4, 6)--(2, 6);
	\node[label] at (0, 0) {\small 0};
	\node[label] at (0.5, 1) {\small 1};
	\node[label] at (1, 2) {\small 2};
	\node[label] at (1.5, 3) {\small 3};
	\node[label] at (2, 4) {\small 4};
	\node[label] at (2.5, 5) {\small 5};
	\node[label] at (3, 6) {\small 6};
	\end{tikzpicture}
\end{center}
તો આ વર્ણન શા માટે સ્વીકારવું?

એક કારણ એ છે કે તે સરસ અને સહેલાઈથી કામમાં લઈ શકાય એવું વર્ણન
છે. સૌથી સ્વાભાવિક લાગતા વર્ણન, એટલે સહજ ગણરચના,નું પતન થયા
પછી આપણને કોઈ સરસ વર્ણનની જરૂર છે.

પરંતુ આ વર્ણન \emph{માત્ર} સરસ નથી. તેના પર આધારિત કોઈ પણ
ગણસિદ્ધાંત \emph{સુસંગત} હશે એમ માનવા માટે આપણી પાસે સારું
કારણ છે. તેનું કારણ આ છે.

ગણની સંચયી-પુનરાવર્તી કલ્પના મુજબ આપણે ગણો તબક્કાઓમાં રચીએ
છીએ; અને તેમના !!{element}s \emph{પહેલેથી} ઉપલબ્ધ વસ્તુઓ
હોવા જોઈએ. તેથી કોઈ પણ તબક્કા~$S$ માટે આ ગણ રચી શકીએ:
\[
R_S = \Setabs{x}{x \notin x \text{ અને $x$ તબક્કા $S$ પહેલાં ઉપલબ્ધ હતો}}
\]
રસેલનો વિરોધાભાસ સાબિત કરવામાં વપરાયેલું તર્ક હવે સ્થાપિત
કરશે કે $R_S$ પોતે તબક્કા $S$ પહેલાં ઉપલબ્ધ નથી. અને તે
પરસ્પરવિરોધ નથી. વધુમાં, ગણની સંચયી-પુનરાવર્તી કલ્પના
સ્વીકારીએ તો રસેલનો ગણ પોતે રચી શકવાની \emph{અપેક્ષા} પણ
ન રાખવી જોઈએ. કારણ કે તે ``ક્યારેય પણ ઉપલબ્ધ થવાના'' બધા
સ્વસભ્ય ન હોય એવા ગણોનો ગણ હોત. ટૂંકમાં, રસેલનો ગણ આપણે
(સાબિતી મુજબ) રચી શકતા નથી તે સંચયી-પુનરાવર્તી વર્ણન મુજબ
\emph{આશ્ચર્યજનક} નથી; એ તો આપણે કરેલી \emph{આગાહી} જ છે.

%In one sense, then, the cumulative-iterative conception of set yields a response to Russell's Paradox which is rather like the \emph{predicativist}'s response. After all, the predicativist said that the collection of all non-self-membered sets$_n$ is a set$_{n+1}$, and this set$_n$ will not be self-membered. (Indeed, most predicativists will treat it as \emph{ungrammatical} to try to ask whether a set is self-membered.)
%
%But there is an important difference between the cumulative-iterative approach and the predicativist approach: the cumulative-iterative approach treats all of our entities as being \emph{of the same kind}. The predicativist will presumably have an empty set$_0$ (i.e., a set$_0$ with no !!{element}s), and an empty set$_1$ (i.e., a set$_1$ with no !!{element}s), and these will be \emph{different entities}. But on the cumulative-iterative approach, there is just one empty set, $\emptyset$, and it will be available at every stage of the hierarchy. 

%Indeed, the Axiom of Extensionality, stated at the start of this chapter, will hold true of the elements in this hierarchy (so that there can be \emph{only one} ``empty'' set). But I do not intend to use the cumulative-iterative conception to justify Extensionality (see \cite{Boolos1971}). Again: I suggest that we should accept Extensionality, just because we are interested in extensional collections. 

\end{document}
